% !TEX root = ../main.tex
\section{Experimental Landscape}
\label{sec:landscape}

% Spine (2026-08-09 grilling; next-gen/close 2026-08-27): solar/reactor -> LBL
% (NOvA/T2K; inline T2K--NOvA joint + JUNO--LBL |dm^2| ordering hook; brief
% IceCube) -> DUNE-weighted next-generation (DUNE Phase I/II + exposure-indexed
% reach; Hyper-K short; complementarity) -> NuFIT 6.1 table + pointer
% opens / meantime NOvA disappearance handoff. Intro: provenance claim; body:
% programme sectors. Citation convention: page numbers refer to the arXiv
% version named in each bib entry, not to journal pagination. Sources cited by
% table or figure number use pagination-independent identifiers. See the header
% comment in main.bib. See ADR-0003.

The physics describing neutrino oscillations was set out in the previous section. This
included the three-flavour oscillation framework, which helped to introduce 
the PMNS mixing matrix and the six oscillation parameters that govern flavour change in vacuum.
Matter effects were also introduced and shown to modify the oscillation probabilities. 
This section turns to the current status of our knowledge of these parameters and the 
experiments that constrain them, past, current and future. Many experiments world wide, probing 
the different sectors, which are definted by their distinct $L/E$ ratios and the flavours they can tag,
are each sensitive to their own set of oscillation parameters. This section outlines the
experimental programmes that measure the parameters of each sector, and subsequently 
summarising the global picture. Whilst the discussion will show the current understanding 
which includes precise measurements, three main open questions still remain. Namely, the mass ordering, 
the octant of $\theta_{23}$ and the value of $\dcp$~\cite{NuFIT60,NuFIT61}. Future next genration experiments seek to 
address these open questions, and will also be discussed.
 

% ============================================================
\subsection{Solar and Reactor Measurements}
\label{sec:landscape-solar-reactor}



\paragraph{The solar sector.}
As touched on in Sec.~\ref{sec:history-anomalies}, the solar deficit was first observed
by the Homestake experiment, providing a first indication of neutrino oscillations.
Super-Kamiokande and SNO then provided further compelling experimental evidence for such
flavour change~\cite{Fukuda2001,Ahmad2002}. However, it was only a year later, in 2003, when
the Kamioka Liquid Scintillator Antineutrino Detector (KamLAND) first provided initial
constraints on the solar parameter pair $(\theta_{12},\dm{21})$~\cite{KamLAND2003}. 
Each of these experiments probe the solar sector; with neutrinos produced in nuclear reactions 
in the Sun being the primary source for Super-Kamiokande and SNO. And reactor neutrinos being the 
source for KamLAND. Super-Kamiokande measures neutrinos via $\nu$--electron elastic scattering in 
a large water Cherenkov detector, whilst SNO uses a heavy-water Cherenkov detector to allow for NC, 
CC and elastic scattering interactions to be separately tagged. KamLAND instead measures the disappearance
of reactor $\nuebar$ with a liquid-scintillator detector at a flux weighted average distance of 
$180$\,km~\cite{KamLAND2003}. Super-Kamiokande and SNO therefore probe the solar parameters 
with matter effects via the traversing of the neutrinos through the falling electron density in 
the Sun's core, whilst KamLAND's provides a complementary determination of these parameters 
without such effects. Solar matter effects helped fix the sign of $\dm{21}$ and resolve the octant
of $\theta_{12}$, which reactor experiments like KamLAND could not do alone. $\dm{21}$ was therefore
found to be a positive quantity. Subsequently, KamLAND was able to confirm the values of these
parameters with higher precision in the vacuum case~\cite{KamLAND2013}.





\paragraph{JUNO and the solar sector today.}
\label{sec:landscape-juno}
With the start of data taking in 2025, the first results from the Jiangmen 
Underground Neutrino Observatory (JUNO) have substantially improved the 
solar picture. JUNO is a $20$\,kt liquid-scintillator detector located such that is is 
equidistant to all reactor cores from the Yangjiang and Taishan reactor complexes~\cite{JUNO2026}. 
The reactor $\nuebar$ are detected via inverse beta decay and this baseline was chosen
such that the $L/E$ sits close to the first disappearance minimum.
Figure~\ref{fig:juno-pee-baseline} shows the slow solar and fast atmospheric oscillations in
the $\nuebar$ survival probability as a function of a baseline, considering a fixed 
$4$\,MeV recator neutrino energy. JUNO, DayaBay and KamLAND are placed on this plot for 
comparison.

\begin{figure}[tb]
    \centering
    \includegraphics[width=0.85\textwidth]{slow_fast_nue_osc_2.png}
    \caption{$\nuebar$ survival probability as a function of baseline at fixed $4$\,MeV reactor 
    energy. The red curve indicates the slow oscillations driven by the solar splitting, $\dm{21}$, 
    while the blue curve shows the full survival probability, which includes the fast atmospheric 
    oscillations due to $\dm{31}$. The JUNO baseline of $52.5$\,km sits near the first solar disappearance
    minimum; Daya Bay (near/far) and KamLAND are shown for comparison. Reproduced from the
    JUNO Collaboration~\citep[Fig.~5]{JUNO2026}.}
    \label{fig:juno-pee-baseline}
\end{figure}
After just $59.1$\,days of data collected, JUNO's first results were relased, reporting a
simultaneous high-precision determination of both solar oscillation parameters, $\sin^{2}\theta_{12}$ 
and $\dm{21}$~\cite{JUNO2026}:
\begin{equation}
    \sin^{2}\theta_{12} = 0.3092 \pm 0.0087,
    \qquad
    \dm{21} = (7.50 \pm 0.12)\times 10^{-5}~\text{eV}^{2}.
    \label{eq:juno-solar}
\end{equation}
This result represents a substancial improvement on the precision of the previous combined 
measurements by a factor of 1.6 when considering the normal mass ordering scenario.
Figure~\ref{fig:juno-solar} shows JUNO's $\dm21$ vs $/sin^2 \theta_{12}$ confidence intervals alongside 
results from the previously mentioned experiments. Comparing the error bars for each result shows
that a reactor measurement now leads the measurement of the solar pair of parameters.
However, results from solar data are still very important for the improvement in the determination
of the sign of $\dm{21}$ and the octant of $\theta_{12}$ via the matter effects from neutrino propagation 
in the Sun. 
\begin{figure}[tb]
    \centering
    \includegraphics[trim=302bp 542bp 7bp 64bp,clip,width=0.70\textwidth]{juno_solar_src.pdf}
    \caption{Confidence intervals on $\sin^{2}\theta_{12}$ and $\dm{21}$ from the first
    JUNO data set, with the KamLAND and combined Super-Kamiokande\,+\,SNO determinations
    shown for comparison. The side panels give the one-dimensional $\Delta\chi^{2}$
    profiles. Reproduced from the JUNO Collaboration~\citep[Fig.~4, p.~14]{JUNO2026}.}
    \label{fig:juno-solar}
\end{figure}
It's useful to note that for JUNO's first results, the $\theta_{13}$ and $\dm{31}$ parameters are input 
as external constraints and are not fitted outputs. 
% The Daya Bay short-baseline reactor experiment, provides the leading $\theta_{13}$ measurement, which is what is used in JUNO's fit. 
JUNO's position along the blue curve of Figure~\ref{fig:juno-pee-baseline} is such that it can
also resolve these fast oscillations and hence measure the atmospheric mass splitting $|\dm{31}|$. 
This is a measurement JUNO has performed (\textcolor{red}{no arxiv for $\dm{31}$ result yet}).
Such measurements are also used as external constraints in Long Baseline (LBL) accelerator experiments, 
which will be discussed later.

Super-Kamiokande were the first to establish the parameters pertaining to the atmospheric sector, namely
$\theta_{23}$ and $\dm{32}$. This was achieved by using the zenith angle of atmospheric neutrino 
interactions to estimate their path lengths through the Earth, from down-going to up-going, thereby 
achieving different baselines. Today however, accelerator long-baseline experiments lead the precision 
on these parameters (Sec.~\ref{sec:landscape-lbl}). Since the leading order term in the appearance probability
goes as $P_{\rm atm}$~\hyperref[eq:patm-psol-red]{(\ref{eq:patm-psol-red})} and LBL experiments have weak 
$\theta_{31}$ sensitivity, like JUNO, external constraints on this parameter are essential.

Such a constraint comes from the short-baseline reactor experiment of Daya Bay, in China, which is located
between the Daya Bay and Ling Ao reactor complexes and measures $\nuebar$ disappearance.
At a baseline of $\approx 1.6 km$, the leading order $\dm{31}$ oscillations term sits near
the first disappearance minimum, with the depth of the dip isolating $\theta_{13}$. With two near 
detector halls to constrain the flux and $\nuebar$ spectrum and the far detector halls measuring the $\nuebar$ 
deficits\footnote{These are the two Daya Bay (near/far) markers in Fig.~\ref{fig:juno-pee-baseline}}, 
Daya Bay's result from its final dataset collected over $3158$ days \cite{DayaBay2023} leads the measurement 
for $\theta_{13}$:
\begin{equation}
    \sin^{2}2\theta_{13} = 0.0851 \pm 0.0024.
    \label{eq:dayabay-theta13}
\end{equation}

% Its importance for long-baseline physics is indirect but decisive. Because reactor
% $\nuebar$ disappearance has no $\dcp$ dependence and a negligible matter effect,
% Eq.~\eqref{eq:dayabay-theta13} can be imported into accelerator appearance fits as an
% external prior. Without it the $\nue$ appearance rate would be a function of two
% unknowns at once ($\theta_{13}$ and $\dcp$, entangled also with the mass ordering), and
% neither $\dcp$ nor the ordering could be cleanly extracted. This is why NOvA's
% three-flavour oscillation analyses, discussed in Sec.~\ref{sec:landscape-lbl}, are
% quoted both with and without the Daya Bay $\theta_{13}$ constraint applied.

% ============================================================
\subsection{Long-Baseline Accelerator Measurements}
\label{sec:landscape-lbl}

{\color{red}
\textit{Rough plan (grilling contract):}
\begin{itemize}
    \item NOvA: $\dm{32}$, $\theta_{23}$, weak ordering / $\dcp$ preference
    \item T2K: same parameters, shorter baseline (smaller matter effect)
    \item T2K--NOvA joint: brief \emph{inline} paragraph (no peer joint-fit subsection)
    \item JUNO--LBL $|\dm{3\ell}|$ ordering hook: conceptual only; quantitative later
    \item IceCube DeepCore: one or two sentences (wider programme; not next-gen)
\end{itemize}
}
Unlike solar or reactor experiments, long baseline accelerator experiments have far more control over both 
the baseline and the neutrino energy spectrum. They choose these such that the $L/E$ samples the first 
atmospheric oscillation maximum. The dip in the disappearance spectrum provides the principle 
$\dm{32}$ and $\sin^2(2\theta)$ sensitivity, while the appearance sample allows for the $\dcp$, the ordering and
the octant of $\theta_{23}$ to be inferred.




\paragraph{NOvA.}
The NOvA experiment, whose setup will be discussed in the next chapter, utilises a baseline of $810$\,km
and has an energy spectrum that peaks between the $1$--$3$\,GeV range, such that it sits near that first 
oscillation maximum.
NOvA's latest full 3-flavour oscillation result analyses its full ten-year data set, with an exposure of
$26.61\times10^{20}$ and $12.50\times10^{20}$ protons on target (POT) in neutrino and
antineutrino beam modes respectively. It fits the $\numu$, $\numubar$, $\nue$ and
$\nuebar$ samples simultaneously, the details of which will be discussed in a later chapter. 
NOvA boasts the most precise single experiment constraint on the atmospheric mass squared splitting 
$\dm{32} = 2.431^{+0.036}_{-0.034} (-2.479^{+0.036}_{-0.036})\times 10^{-3}~\text{eV}^{2}$
for the normal (inverted) mass ordering hypothesis~\citep[p.~1]{NOvA2025}.
In both ordering hypotheses, a value for $\sin^{2}\theta_{23} = 0.55^{+0.02}_{-0.02}$, which sits near the 
maximal mixing case, was measured. No significant preference for the octant of $\theta_{23}$ was found.
The normal ordering hypothesis is slightly favoured with a Bayes factor or 2.4 compared to the inverted ordering,
indicating that the NO is 2.4 times more probable than the IO~\citep[p.~1]{NOvA2025}. Applying Daya Bay's
$\theta_{13}$ measurement as a 1-dimensionsl external constraint increases the preference to $3.3$ and adding their
constraint on the atmospheric splitting further raises it to $6.6$. These Bayes Factors are not yet large enough
to consider the ordering resolved.

Figure~\ref{fig:nova-contour} shows the allowed region in the $\dm{32}$--$\sin^{2}\theta_{23}$ space for NOvA,
as well as other experiments, showing NOvA's best fit point, which lies in the region favoured by the normal
ordering hypothesis. The summarised results above can be seen from this plot. The width of NOvA's 
allowed $90 \%$ confidence interval in $\dm{32}$ is the most narrow of all other plotted experiments, reinforcing
its claim as the most precise measurement for this parameter. In contrast, the $90 \%$ confidence region sits 
almost symmetrically about the $\sin^{2}\theta_{23} = 0.5$ line, which indicates maximal mixing, and that the 
octant of $\theta_{23}$, as mentioned, is not well constrained.


\begin{figure}[tb]
    \centering
    \includegraphics[trim=305bp 572bp 27bp 45bp,clip,width=0.66\textwidth]{nova_contour_src.pdf}
    \caption{NOvA's $90\%$ confidence interval in $\dm{32}$--$\sin^{2}\theta_{23}$ space. Plot compares NOvA and
    other experiments' results. NOvA's latest results region is shown with 1D Daya Bay $\theta_{13}$ external 
    constraint. Plot obtained from~\citep[Fig.~6]{NOvA2025}.}
    \label{fig:nova-contour}
\end{figure}

\paragraph{NOvA-T2K.}
The Tokai to Kamioka (T2K) experiment, measures the same parameters as NOvA, as it sits at a similar $L/E$.
T2K has a shorter baseline of $295$\,km and its energy spectrum peaks lower at around $0.6$\,GeV.
Both experiments sit near the first disappearance dip and as a result are broadly consistent in their 
disappearance measurements as can be seen in Figure~\ref{fig:nova-contour}.
Since NOvA's baseline is around $2.7\times$ larger than T2K's, NOvA's neutrinos pass through more of the Earth. The 
appearance channel for NOvA therefore sees an enhanced matter effect. As a consequence, the ordering--$\dcp$ 
degeneracy introduced in Section~\ref{sec:matter} manifests differently in both experiments. The bi-event plot
of Figure~\ref{fig:bievent} shows this more concretely. This plot shows the expected number of $\nuebar$ against $\nue$ 
appearance candidate events at NOvA (a) and T2K's (b) far detectors. By varying the value of $\dcp$ through a full 
period, from $-\pi$ to $\pi$, the predicted number of events traces out an ellipse in this space, with each ordering tracing 
out its own ellipse. The degeneracy explained in Section~\ref{sec:matter}, is clearly visible in the NOvA panel (a), that is, 
the appearance rate for the normal ordering hypothesis with $\dcp\simeq+\pi/2$ yields an almost identical rate to the 
inverted ordering hypothesis with $\dcp\simeq-\pi/2$. For these two hypotheses, the degeneracy occurs due to the matter term 
and the CP term cancelling, resulting in an almost identical appearance rate. With T2K's smaller matter effect, the ordering
sensitivity is weakened, which can be seen by the ellipses overlapping much more than in NOvA. This additional overlap means
that for maximal CP violation, the degeneracy is lifted slightly; the same $\dcp$ values coincide more closely in both 
orderings. For CP-conserving points, the opposite is true; the degeneracy between different CP-conserving values and the 
ordering is still present.

\begin{figure}[tb]
    \centering
    \includegraphics[trim=35bp 424bp 322bp 45bp,clip,width=0.48\textwidth]{novat2k_bievent_src.pdf}
    \caption{``Bi-event'' plot, which shows the predicted number of $\nue$ vs $\nuebar$ appearance
    candidate events for NOvA (a) and T2K (b). Varying $\dcp$ through its full range traces an ellipse,
    with four reference values of $\dcp$ marked, with one ellipse for the normal (blue), and one
    for the inverted (orange) orderings. Taken from the T2K and NOvA Collaborations~\citep[Fig.~1, p.~6]{T2KNOvA2025}.}
    \label{fig:bievent}
\end{figure}

Comparing the data points for both experiments, NOvA's point is closest to the overlap of the ellipses, occupying the most 
degenerate region, whereas T2K's is placed at the tip of the normal ordering ellipse corresponding to $\dcp\simeq-\pi/2$ in 
either ordering. Due to the pulling apart of the ellipses, NOvA is more sensitive to the ordering, whereas T2K's overlap makes
its $\dcp$ sensitivity stronger. For this reason, the joint analysis of the two experiments combines the appearance samples.


This first joint analysis yielded a value for $\dm{32} = 2.43^{+0.04}_{-0.03}(-2.48^{+0.03}_{-0.04})\times10^{-3}$\,eV$^{2}$ in 
the normal (inverted)ordering, with a $3\sigma$ interval on $\dcp$ of $[-1.38\pi,\,0.30\pi]$ in the normal ordering, narrowing 
to $[-0.92\pi,\,-0.04\pi]$ when assuming the inverted ordering ~\citep[p.~5]{T2KNOvA2025}. The joint collaboration performed a 
joint Bayesian posterior fit over a shared set of oscillation parameters. The fit showed no significant preference for either 
ordering. However, when the inverted ordering hypothesis is assumed, the fit excludes a substantial fraction of $\dcp$ space, 
and favours maximal CP violation with $\dcp\simeq-\pi/2$~\citep[Fig.~3, p.~9]{T2KNOvA2025}.


\paragraph{Reactor--accelerator complementarity.}
The determination of the atmospheric mass squared splitting parameter $\dm{3\ell}$ by reactor and long baseline
accelerator experiments are complementary in that their measurements should agree under the correct mass
ordering. Reactor $\nuebar$ disappearance experiments are sensitive to $|\dm{31}|$, whereas long-baseline $\numu$
disappearance measures $\dm{32}$. These two parameters, arising as frequencies in the oscillation probability, differ 
by the known solar splitting $\dm{21}$. Therefore, the two types of experiments should agree after subtracting the solar splitting
from $\dm{31}$ under the correct mass ordering hypothesis. A reactor measurement of $\dm{31}$ can therefore be used to help with 
the mass ordering determination in NOvA, an idea that will be fleshed out in a later chapter. As a results, as reactor experiments 
like JUNO matures and becomes more precise, NOvA will be able to improve its ordering preference~\cite{NOvAJUNO2026}.

% ============================================================
\subsection{Next-Generation Experiments}
\label{sec:landscape-nextgen}

\paragraph{DUNE.}
The Deep Underground Neutrino Experiment (DUNE) will send a wide-band neutrino beam $1300$\,km from Fermilab to
a far detector complex consisting of up to four liquid-argon time-projection chambers located at the Sanford
Underground Research Facility, with a movable near detector (DUNE-PRISM) on the Fermilab site~\cite{DUNE2025Science}. Matter effects are enhanced due to the much longer baseline
when compared to NOvA and T2K. DUNE's bi-event plot would show the ellipses no longer overlapping, meaning that no pair
of $\dcp$ values would coincide in both orderings. The separation of the ellipses indicates DUNE's highly enhanced
ordering sensitivity. In contrast, DUNE's $\dcp$ sensitivity cannot be enhanced enough by matter effects alone.
Even after the ellipses have been fuly separated thanks to increased matter effects, a single data point, lying on an ellipse 
will still have some degeneracy (in $\dcp$) caused by a change in $\theta_{23}$. NOvA and T2K both sample a single oscillation maximum, where the ordering-dependent matter and the $\dcp$ interference terms can appear to cancel (Sec.~\ref{sec:matter}). To overcome these degeneracies, DUNE will make use of a wide-band beam, covering both the first and second oscillation maxima as well as studying both the appearance and disappearance spectra. The first maximum allows the ordering to be determined via the matter effect, while the second maximum shows an enhanced CP asymmetry (Fig.~\ref{fig:dune-appearance}), through which DUNE's $\dcp$ sensitivity enters.

\begin{figure}[tb]
    \centering
    \includegraphics[trim=46bp 466bp 333bp 166bp,clip,width=0.48\textwidth]{dune_appearance_src.pdf}\hfill
    \includegraphics[trim=46bp 238bp 333bp 394bp,clip,width=0.48\textwidth]{dune_appearance_src.pdf}
    \caption{The $\numu\to\nue$ (left) and $\numubar\to\nuebar$ (right) appearance
    probabilities at DUNE's $1285$\,km baseline in the normal ordering, for
    $\dcp=-\pi/2$ (blue), $0$ (red) and $\pi/2$ (green). The first oscillation
    maximum is the broad peak near $2.5$\,GeV; the second is the sharper peak
    near $0.8$\,GeV. Reproduced from the DUNE Collaboration~\citep[Fig.~1, p.~8]{DUNE2020}.}
    \label{fig:dune-appearance}
\end{figure}

DUNE will also measure $\theta_{23}$ with world-leading precision allowing for the octant to be determined and the wide-band beam will allow for the remaining degeneracy between $\dcp$ and $\theta_{23}$ to be resolved~\citep[p.~3]{DUNE2025Science}. With the complementary $\theta_{23}$ dependence of both oscillation channels (Sec.~\ref{sec:three-flavour}), DUNE’s appearance and disappearance measurements together will probe the octant~\citep[p.~28]{DUNE2020}. To achieve the target level of precision, the neutrino--argon cross section as a function of energy must be determined at the few-percent level. This shared cross-section systematic is constrained by the near detector; the DUNE-PRISM system moves the near LArTPC off axis allowing it to sample a range of off-axis energies spectra~\cite{DUNE2025Science}.

\paragraph{Hyper-Kamiokande.}
Hyper-Kamiokande will share the same baseline as T2K but will use a significantly larger water-Cherenkov detector~\cite{HyperK2025}. 
Like with T2K, since the baseline is short, so are the matter effects. The appearance channel therefore doesn't have high ordering 
sensitivity but will primarily measure $\dcp$. The ordering sensitivity will instead arise from atmospheric neutrino, since they 
will feel an enhanced matter effect due to their distances traveled through the Earth. 

Both next generation experiments seek to measure the same parameters, but in different ways, with different detector technologies, baselines and beam energy spectra. Therefore, agreement between them would point towards a correct shared understanding of neutrino oscillations, while disagreement could point towards issues with one of the experiments, or an incomplete three-flavour picture~\citep[p.~3]{DUNE2025Science}.

% ============================================================
\subsection{Current Constraints and Open Questions}
\label{sec:landscape-open}

Table~\ref{tab:osc-params} collects the current parameter values. The global-fit column
is taken from NuFIT~6.1, which combines data available through its November~2025
closing date (including the first JUNO solar-pair result); the second column gives the
leading single-experiment determination for each parameter.

\begin{table}[tb]
    \centering
    \small
    \begin{tabular}{lccll}
        \toprule
        & \multicolumn{2}{c}{NuFIT~6.1 global fit} & \multicolumn{2}{c}{Leading direct} \\
        \cmidrule(lr){2-3}\cmidrule(lr){4-5}
        Parameter & Normal & Inverted & Value & Source \\
        \midrule
        $\sin^{2}\theta_{12}$ & $0.3088^{+0.0067}_{-0.0066}$ & $0.3088^{+0.0067}_{-0.0066}$
            & $0.3092 \pm 0.0087$ & JUNO \\
        $\sin^{2}\theta_{23}$ & $0.470^{+0.017}_{-0.014}$ & $0.550^{+0.013}_{-0.016}$
            & $0.55 \pm 0.02$ & NOvA \\
        $\sin^{2}\theta_{13}$ & $0.02248^{+0.00055}_{-0.00059}$ & $0.02262^{+0.00057}_{-0.00056}$
            & $0.0851 \pm 0.0024$\rlap{$^{\dagger}$} & Daya Bay \\
        $\dcp/^{\circ}$ & $212^{+26}_{-36}$ & $274^{+22}_{-25}$
            & --- & \\
        $\dm{21}/10^{-5}~\text{eV}^{2}$ & $7.537^{+0.094}_{-0.10}$ & $7.537^{+0.094}_{-0.10}$
            & $7.50 \pm 0.12$ & JUNO \\
        $\dm{3\ell}/10^{-3}~\text{eV}^{2}$ & $+2.511^{+0.021}_{-0.020}$ & $-2.483^{+0.020}_{-0.020}$
            & $+2.431^{+0.036}_{-0.034}$ & NOvA \\
        \bottomrule
    \end{tabular}
    \caption{Three-flavour oscillation parameters. Global-fit values are the NuFIT~6.1
    best-fit points with $1\sigma$ uncertainties for the normal (NO) and inverted (IO)
    orderings, from the variant including Super-Kamiokande and IceCube atmospheric
    data~\cite{NuFIT61,NuFIT60}; that analysis closes in November~2025. The final
    columns give the leading single-experiment determination:
    JUNO~\citep[p.~15]{JUNO2026}, NOvA~\citep[p.~2]{NOvA2025} and
    Daya Bay~\cite{DayaBay2023}. Here $\dm{3\ell}$ denotes $\dm{31}$ for the normal and
    $\dm{32}$ for the inverted ordering. $^{\dagger}$Daya Bay reports
    $\sin^{2}2\theta_{13}$ rather than $\sin^{2}\theta_{13}$.}
    \label{tab:osc-params}
\end{table}

Three questions remain open.

\paragraph{The mass ordering.}
The sign of $\dm{32}$ is still open. Fig.~\ref{fig:bievent} shows why a single baseline
cannot separate it from $\dcp$, and the joint T2K--NOvA analysis finds no strong
preference either way~\citep[p.~5]{T2KNOvA2025}.

\paragraph{The $\theta_{23}$ octant.}
Disappearance data constrain $\sin^{2}(2\theta_{23})$ and are therefore blind to the
octant at leading order. Fig.~\ref{fig:nova-contour} shows the present status: NOvA's
allowed region sits almost symmetrically about maximal mixing, which is the hardest case
to resolve.

\paragraph{$\dcp$.}
The value of $\dcp$ is the least well determined of the six parameters. Its allowed range
depends on the assumed ordering, as the joint T2K--NOvA interval already showed, and
Fig.~\ref{fig:bievent} shows why. CP conservation is not yet excluded.

Until DUNE and Hyper-Kamiokande resolve these questions, they must be constrained by the
experiments now running. NOvA's disappearance measurement is the leading direct
determination of $\dm{32}$~\citep[p.~1]{NOvA2025}, and it is that measurement --- the
far-detector spectrum, its systematic treatment and its interpretation --- that the
remainder of this thesis addresses.
